\lectures{9,10,11}
\title{The Maximum Likelihood Estimator}

\begin{document}

\subsection{Types of Parameters, Identifiability}

\textbf{Definition.} In a statistical model,
the parameters we care about are called \textbf{parameters of interest},
and those we don't are called \textbf{nuisance parameters}.
A parameter $\theta$ is \textbf{identifiable} if
distinct values of $\theta$ necessarily yield distinct probability distributions.

\textbf{Example.} In the model
$\{\mathcal{N}(\mu, \sigma^2) \mid \mu \in \mathbb{R}, \sigma \in \mathbb{R}^+\}$,
we might say $\mu$ is a parameter of interest and $\sigma$ is a nuisance parameter.

\textbf{Example.} The response to a drug dosage is modeled by
$\mathcal{N}\left(E_0 + \frac{D \times E_{\text{max}}}{D + ED_{50}},\ \sigma^2\right)$,
for some parameter $\theta = (\sigma^2, E_0, E_{\text{max}}, ED_{50})$.
In this case, $\theta$ is not identifiable.

\subsection{Defining the Maximum Likelihood Estimator (MLE)}

Consider some (i.i.d.) random variable samples $X_1, \dots, X_n \sim \mathbb{P}$
and a parameter $\theta$.  Furthermore, some new notation:

\textbf{Notation.} Given a parameter $\theta$,
let its true, population value be denoted by $\theta^*$.

If the parameter $\theta$ can be expressed using $\mathbb{E}[\dots]$ only,
it's not hard to define an okay estimator.

\[ \text{(ex:)} ~~~~~~ \theta := \frac{\mu}{\sigma} =
\frac{\mathbb{E}[X_i]}{(\mathbb{E}[X_i^2] - \mathbb{E}[X_i]^2)^{1/2}} ~ \implies ~ \hat{\theta}_n :=
\frac{\frac{1}{n}\sum_{i = 1}^n X_i}{\left(
\frac{1}{n}\sum_{i = 1}^n X_i^2 - \left(\frac{1}{n}\sum_{i = 1}^n X_i\right)^2
\right)^{1/2}}. \]

More smartly, though, we should think to use the \textbf{maximum likelihood estimator (MLE)}.

\textbf{Definition.} Let $f_{\theta}(x)$ denote the PDF corresponding to
the distribution $\mathbb{P}_{\theta}$, and
consider some samples $X_1, \dots, X_n \sim \mathbb{P}_{\theta^*}$.
\begin{itemize}
    \item The \textbf{likelihood} function is defined by
    $L_n(\theta) := \prod_{i = 1}^n f_{\theta}(X_i)$.
    \item The \textbf{log likelihood} function is just $\ell_n(\theta) := \log L_n(\theta)
    = \sum_{i = 1}^n \log f_{\theta}(X_i)$.
    \item The \textbf{maximum likelihood estimator (MLE)} $\hat{\theta}_n$ is the value
    of $\theta$ for which $L_n(\theta)$ (or $\ell_n(\theta)$) is maximized.
\end{itemize}
In other words, the MLE $\hat{\theta}_n$
is the value for the parameter $\theta$ whose distribution $\mathbb{P}_{\hat{\theta}_n}$
is most ``compatible'' with the samples taken from $\mathbb{P}_{\theta^*}$,
which would suggest that $\mathbb{P}_{\hat{\theta}_n}$
might be a good representation of $\mathbb{P}_{\theta^*}$.

\begin{remark}
    Reason through the following and check that they make sense.
    \begin{itemize}
        \item For the model $\{\mathrm{Ber}(\theta) \mid \theta \in [0, 1]\}$,
        the MLE is $\hat{\theta} := \frac{1}{n} \sum_{i = 1}^n X_i$.
        \item For the model $\{\mathcal{N}(\theta, 1) \mid \theta \in \mathbb{R}\}$,
        the MLE is $\hat{\theta} := \frac{1}{n} \sum_{i = 1}^n X_i$.
        \item For the model $\{\mathrm{Unif}([0, \theta]) \mid \theta \geq 0\}$,
        the MLE is $\hat{\theta} := \max X_i$.
    \end{itemize}
\end{remark}

Note that $L_n(\theta)$ and $\ell_n(\theta)$ are
\emph{deterministic} functions whose behaviors are determined by \emph{random} samples
from $\mathbb{P}_{\theta^*}$.
If one seeks a likelihood function that is instead determined by $\theta^*$ itself,
one may opt for\dots

\textbf{Definition.} The \textbf{population log likelihood} function is 
 $\ell(\theta) := \mathbb{E}_{\theta^*}[\log f_{\theta}(X)]$.
(The $\mathbb{E}_{\theta^*}$ notation
indicates that $X \sim \mathbb{P}_{\theta^*}$.)

\subsection{Kullback-Leibler (KL) Divergence and MLE Consistency}

Rather than maximizing the likelihood, the MLE can also be defined
as the $\theta$ that minimizes the ``distance'' between
$\mathbb{P}_{\theta^*}$ and $\mathbb{P}_{\hat{\theta}}$.
Here's how we formalize that notion of ``distance''.

\textbf{Definition.} Given distributions $\mathbb{P}_1$ and $\mathbb{P}_2$
with PDFs $f_1(x)$ and $f_2(x)$, their \textbf{Kullback-Leibler (KL) divergence} is:
\begin{align*} D_{\mathrm{KL}}(\mathbb{P}_1 \parallel \mathbb{P}_2) :=
\mathbb{E}_{f_1}\left[ \log\left(\frac{f_1(X)}{f_2(X)}\right)\right] & =
\int_{\mathcal{R}} f_1(x) \log \left(\frac{f_1(x)}{f_2(x)}\right) \, \mathrm{d}x \\ & =
\left( \int_{\mathcal{R}} f_1(x) \log f_1(x) \, \mathrm{d}x \right) - 
\left( \int_{\mathcal{R}} f_1(x) \log f_2(x) \, \mathrm{d}x \right). \end{align*}
Here, $\mathcal{R}$ is the entire support of $\mathbb{P}_1$,
which is also the entire support of $\mathbb{P}_2$.
Also, as usual, the $\mathbb{E}_{f_1}$ notation
indicates that the random variable $X$ is to be sampled using the PDF $f_1$.

Note that $D_{\mathrm{KL}}$ is not a metric; it is not symmetric, nor does it
satisfy the triangle inequality.  But Jensen's inequality implies it is at least nonnegative,
with $D_{\mathrm{KL}}(\mathbb{P}_1 \parallel \mathbb{P}_2) = 0$ if and only if
$\mathbb{P}_1 = \mathbb{P}_2$.

\begin{remark}
    Jensen's inequality states that
    for any concave function $\varphi$ and any random variable $Z$,
    we have $\mathbb{E}[\varphi(Z)] \leq \varphi(\mathbb{E}[Z])$.
    Equality is precisely when $Z$ is almost-surely constant,
    assuming $\varphi$ is \emph{strictly} concave.
\end{remark}

Importantly, this implies the following.

\begin{theorem}{KL Divergence Minimization}
    The map $\theta \mapsto
    D_{\mathrm{KL}}(\mathbb{P}_{\theta^*} \parallel \mathbb{P}_{\theta})$
    is minimized when $\theta = \theta^*$.
    Also, the population log likelihood
    $\ell(\theta) := \mathbb{E}_{\theta^*}[\log f_{\theta}(X)]$
    is maximized at $\theta = \theta^*$.
\end{theorem}

\begin{proof}
    The first sentence is true by Jensen's; the point is that $\log$ is concave, so:
    \[ D_{\mathrm{KL}}(\mathbb{P}_{\theta^*} \parallel \mathbb{P}_{\theta})
    = -\mathbb{E}_{\theta^*} \left[ \log \left(\frac{f_{\theta}(X)}{f_{\theta^*}(X)}\right) \right]
    \geq -\log \left( \mathbb{E}_{\theta^*} \left[ \frac{f_{\theta}(X)}{f_{\theta^*}(X)} \right] \right)
    = -\log(1) = 0. \]
    Equality occurs precisely when $\frac{f_{\theta}(X)}{f_{\theta^*}(X)}$
    (with $X \sim \mathbb{P}_{\theta^*}$) is almost-surely constant,
    or when $f_{\theta} \cong f_{\theta^*}$, as desired.

    The second sentence follows because minimizing
    $D_{\mathrm{KL}}(\mathbb{P}_{\theta^*} \parallel \mathbb{P}_{\theta})$
    is equivalent to maximizing
    $\int_{\mathbb{R}} f_{\theta^*}(x) \log f_{\theta}(x) \, \mathrm{d}x$;
    see the far-RHS of the definition of $D_{\mathrm{KL}}$.
    But $\int_{\mathbb{R}} f_{\theta^*}(x) \log f_{\theta}(x) \, \mathrm{d}x = \ell(\theta)$
    identically. ~ $\blacksquare$
\end{proof}

The upshot of introducing KL Divergence is the second sentence of the
``KL Divergence Minimization'' theorem.
We will leverage this fact to prove that the MLE is always a consistent estimator.

\begin{theorem}{MLE Consistency}
    The MLE is consistent; that is, if $\hat{\theta}_n$ is the MLE of the parameter $\theta$,
    then $\hat{\theta}_n \convprob \theta^*$ always.
\end{theorem}

\begin{proof}
    Recall the following.
    \begin{itemize}
        \item The MLE $\hat{\theta}_n$ is
        the maximizer of the log likelihood $\ell_n(\theta)$ by definition.
        \item The true parameter value $\theta^*$ is
        the maximizer of the population log likelihood $\ell(\theta)$
        by the previous theorem.
        \item For \emph{any} value $\theta'$,
        by the LLN, the sequence $\{\frac{1}{n}\ell_n(\theta')\}_{n = 1}^{\infty}$
        converges in probability to the constant $\ell(\theta')$.
    \end{itemize}
    These three points together imply $\hat{\theta}_n \convprob \theta^*$.
    Roughly speaking, the third bullet point says $\ell_n \convprob \ell$,
    so the maximizer of $\ell_n$ should also converge in probability
    to the maximizer of $\ell$.
    (Yes, this needs to be justified---handwaving a little here.) ~ $\blacksquare$
\end{proof}

\subsection{Fisher Information}

We claim the MLE also achieves asymptotic normality.  In fact,
we claim we even know its asymptotic variance.
It turns out, however, that we'll need a little more setup before we get to asymptotic variance.

\textbf{Definition.} Consider a \underline{statistical model} with a parameter $\theta$
and varying PDFs given by $f_{\theta}(x)$, dependent on $\theta$.
Given an observation $x$ and a proposed value $\theta$ of the parameter,
the \textbf{score} is:
\[ s(x; \theta) := \frac{\partial}{\partial\theta} \log f_{\theta}(x)
= \frac{\frac{\partial}{\partial\theta} f_{\theta}(x)}{f_{\theta}(x)}. \]
The RHS comes from a single application of the chain rule.
When $s(x; \theta)$ is positive, it means that
$x$ would be better explained if the proposed $\theta$ were increased, and vice versa.

\begin{remark}
    The notation is slightly confusing;
    as a simpler example, consider defining $g(x) := \frac{d}{dx} x^2$ and concluding $g(5) = 10$.
\end{remark}

\begin{theorem}{Zero Mean Score}
    We have $\mathbb{E}_{\theta^*}[s(X; \theta^*)] = 0$,
    where the $\mathbb{E}_{\theta^*}$ notation indicates that
    $X$ is sampled from the distribution $\mathbb{P}_{\theta^*}$;
    in other words, the proposed value $\theta^*$ is correct.

    In this case, the score has some positives and negatives, but expectedly is overall neutral
    in how the proposal $\theta^*$ should change.
\end{theorem}

\begin{proof}
    Just some integral manipulation.
    \[ \mathbb{E}_{\theta^*}[s(x; \theta^*)]
    = \int_{\mathcal{R}} f_{\theta^*}(x) \cdot
    \left[\frac{\frac{\partial}{\partial\theta} f_{\theta}(x)}{f_{\theta}(x)}\right]_{\theta = \theta^*}
    \mathrm{d}x
    = \int_{\mathcal{R}} \left[\frac{\partial}{\partial\theta} f_{\theta}(x)\right]_{\theta = \theta^*}
    \mathrm{d}x
    = \frac{\partial}{\partial\theta}
    \left[\int_{\mathcal{R}} f_{\theta}(x) \, \mathrm{d}x\right]_{\theta = \theta^*}
    = \frac{\partial}{\partial\theta} [1] = 0. ~~~ \blacksquare \]
\end{proof}

Note that the definition of score doesn't care whether $x$ is sampled
from a distribution that uses the proposed value $\theta$ for its parameter.
The definition of Fisher information, on the other hand, \emph{does} care.

\textbf{Definition.} Consider a \underline{statistical model} with a parameter $\theta$
and varying PDFs given by $f_{\theta}(x)$, dependent on $\theta$.
As a function of the true value of $\theta$, the \textbf{Fisher information}
of the resulting distribution is:
\[ I(\theta) := \mathbb{V}_{\theta}\left[s(X; \theta)\right]
= \mathbb{E}_{\theta}[s(X; \theta)^2]. \]
As usual, the $\mathbb{E}_{\theta}$ and $\mathbb{V}_{\theta}$ notation
indicates that $X \sim \mathbb{P}_{\theta}$; in other words, that $\theta$ is the ``true value''
of the parameter.  It's because of this and the previous theorem
that the above two expressions for $I(\theta)$ are equivalent.

When $I(\theta)$ is large,
it means a single observation $x$ will tend to suggest
very \emph{large} changes to the proposed $\theta$.
In other words, singular observations $x$ carry more information.

\begin{problem}
    Given the model $\{\mathcal{N}(\theta, 9) \mid \theta \in \mathbb{R}\}$,
    determine $I(\theta)$.  Do the same for $\{\mathcal{N}(\theta, 900) \mid \theta \in \mathbb{R}\}$.
\end{problem}

\begin{solution}
    Recall that the PDF of $X \sim \mathcal{N}(\theta, 9)$ (with $\sigma = 3$) is given by:
    \[ f_{\theta}(x) = \frac{1}{3\sqrt{2\pi}} \exp\left(-\frac{(x - \theta)^2}{2 \times 3^2}\right)
    ~ \implies ~ s(x; \theta) = \frac{\partial}{\partial\theta} \log f_{\theta}(x)
    = \frac{x - \theta}{3^2}. \]
    The Fisher information is therefore:
    \[ I(\theta) = \mathbb{V}_{\theta}\left[ \frac{X - \theta}{3^2} \right]
    = \left(\frac{1}{3^2}\right)^2 \mathbb{V}_{\theta}[X - \theta]
    = \left(\frac{1}{3^2}\right)^2(3^2) = \frac{1}{9}. \]
    It's a constant!
    Expectedly, the Fisher information for $\{\mathcal{N}(\theta, 900) \mid \theta \in \mathbb{R}\}$
    is constantly $\frac{1}{900}$.
    
    Notably, the Fisher information of the latter is smaller than the former.
    This should make sense; a single observation from
    $\{\mathcal{N}(\theta, 900) \mid \theta \in \mathbb{R}\}$
    tells you a lot less about $\theta$ than a single observation from
    $\{\mathcal{N}(\theta, 9) \mid \theta \in \mathbb{R}\}$ does.
\end{solution}

The Fisher information has one more alternative representation.

\begin{theorem}{Equivalent Fisher}
    Fisher information equals the negative expected curvature of the log-likelihood; that is,
    \[ I(\theta) = -\mathbb{E}_{\theta}\left[
        \frac{\partial^2}{\partial\theta^2} \log f_{\theta}(X)
    \right]. \]
\end{theorem}

\begin{proof}
    Just a lot of calculus.  By the quotient and chain rules, we have:
    \[ \frac{\partial^2}{\partial\theta^2} \log f_{\theta}(X)
    = \frac{\partial}{\partial\theta} \frac{\frac{\partial}{\partial\theta} f_{\theta}(X)}{f_{\theta}(X)} 
    = \frac{f_{\theta}(X) \frac{\partial^2}{\partial\theta^2} f_{\theta}(X)
    - \left(\frac{\partial}{\partial\theta} f_{\theta}(X)\right)^2}{
        \left(f_{\theta}(X)\right)^2
    }
    = \frac{\frac{\partial^2}{\partial\theta^2} f_{\theta}(X)}{f_{\theta}(X)}
    - s(X; \theta)^2. \]
    The expectation of the second term of the RHS is $-I(\theta)$,
    so it suffices to show the expectation of the first term is zero.
    \[ \mathbb{E}_{\theta}\left[\frac{\frac{\partial^2}{\partial\theta^2} f_{\theta}(X)}{f_{\theta}(X)}\right]
    = \int_{\mathcal{R}} f_{\theta}(x) \cdot \left[
        \frac{\frac{\partial^2}{\partial\theta^2} f_{\theta}(x)}{f_{\theta}(x)}
    \right] \, \mathrm{d}x
    = \int_{\mathcal{R}} \left[\frac{\partial^2}{\partial\theta^2} f_{\theta}(x)\right] \, \mathrm{d}x
    = \frac{\partial^2}{\partial\theta^2} \left[\int_{\mathcal{R}} f_{\theta}(x) \, \mathrm{d}x\right]
    = \frac{\partial^2}{\partial\theta^2} [1] = 0. \]
    A very similar argument as before, actually\dots ~ $\blacksquare$
\end{proof}

Finally, when we have multiple samples,
we can talk about ``total Fisher information''.

\textbf{Definition.} Given $n$ samples $\{X_i\}_{i = 1}^n$,
recall $\ell_n(\theta) := \sum_{i = 1}^n \log f_{\theta}(X_i)$
is the log likelihood.
Then the \textbf{total Fisher information} gained from these samples is:
\[ I_n(\theta) :=
-\mathbb{E}_{\theta}\left[\frac{\partial^2}{\partial\theta^2}\ell_n(\theta)\right]
= n I(\theta). \]
These two expressions are equal because all the $X_i$ are sampled from
the same distribution $\mathbb{P}_{\theta}$, so applying linearity of expectation
yields $n$ copies of $I(\theta)$ summed together.

\subsection{MLE Asymptotic Normality}

Now we can finally get to asymptotic normality of the MLE.

\begin{theorem}{MLE Asymptotic Normality}
    The MLE $\hat{\theta}_n$ has asymptotic variance $\frac{1}{I(\theta^*)}$; that is,
    \[ \sqrt{n}(\hat{\theta}_n - \theta^*) \rightsquigarrow
    \mathcal{N}\left(0, \frac{1}{I(\theta^*)}\right). \]
\end{theorem}

\begin{proof}
    Recall the log likelihood
    $\ell_n(\theta) := \sum_{i = 1}^n \log f_{\theta}(X_i)$.
    The definition of the MLE $\hat{\theta}_n$ is that it maximizes $\ell_n(\theta)$, meaning:
    \[ \frac{\partial}{\partial\theta} \ell_n(\hat{\theta}_n)
    = \sum_{i = 1}^n s(X_i; \hat{\theta}_n) = 0. \]
    Recall that we proved the consistency of $\hat{\theta}_n$ earlier,
    meaning $\hat{\theta}_n \convprob \theta^*$.
    Thus, it makes sense to take the Taylor approximation:
    \[ 0 = \sum_{i = 1}^n s(X_i; \hat{\theta}_n)
    \approx \sum_{i = 1}^n s(X_i; \theta^*) +
    (\hat{\theta}_n - \theta^*) \left[
        \frac{\partial}{\partial\theta} \sum_{i = 1}^n s(X_i; \theta)
    \right]_{\theta = \theta^*}. \]
    Moving terms around yields the approximation:
    \[ \sqrt{n}(\hat{\theta}_n - \theta^*)
    \approx - \left(\frac{\frac{1}{\sqrt{n}}\cdot \sum_{i = 1}^n s(X_i; \theta^*)}{
        \frac{1}{n} \cdot \sum_{i = 1}^n \left[\frac{\partial^2}{\partial\theta^2} \log f_{\theta}(X_i)
        \right]_{\theta = \theta^*}
    }\right). \]
    All that setup pays off here;
    it turns out the numerator and denominator of this fraction
    can be expressed very nicely.
    \begin{itemize}
        \item By ``Zero Mean Score'', the definition of Fisher information, and CLT,
        the numerator converges ($\rightsquigarrow$) to $\mathcal{N}(0, I(\theta^*))$.
        \item By ``Equivalent Fisher'' and LLN,
        the denominator converges ($\convprob$) to $-I(\theta^*)$.
    \end{itemize}
    Thus, by Slutsky's lemma, we have:
    \[ \sqrt{n}(\hat{\theta}_n - \theta^*) \rightsquigarrow
    -\left(\frac{\mathcal{N}(0, I(\theta^*))}{-I(\theta^*)}\right)
    = \mathcal{N}\left(0, \frac{1}{I(\theta^*)}\right). \]
    And that's it. ~ $\blacksquare$
\end{proof}

Knowing that the asymptotic variance is $\frac{1}{I(\theta^*)}$
is useful for constructing confidence intervals using the MLE $\hat{\theta}_n$,
especially when the variance of $\hat{\theta}$ would not otherwise be easy to compute directly.

\begin{remark}
    For the model $\{\mathrm{Ber}(p) \mid p \in (0, 1)\}$,
    the Fisher information is $I(p) = \frac{1}{p(1 - p)}$.
    CYU: Why is this expected?
\end{remark}

\end{document}